\section{SmallWood Programming Model}
\label{sec:smallwood-model}
\label{sec:smallwood}

This section gives the SmallWood formulation of the ARC-SPRUCE NIZKs.  The prover commits to packed witness
rows, represented as low-degree row polynomials over a packing set \(\Omega\) and masked on a pairwise-disjoint
hiding set \(\Omega'\).  The verifier then checks PACS-style constraints on those committed rows.  We record here
the semantic statement layers and the exact coefficient-domain consequences needed by the security arguments;
concrete row inventories, gadget-level realizations, and proof-stack engineering are deferred to
Appendix~\ref{app:smallwood}.

\subsection{Proof interface}
\label{subsec:smallwood-interface}
Fix a packing set \(\Omega=\{\omega_1,\ldots,\omega_s\}\subset\Fq\) and a disjoint hiding set
\(\Omega'=\{\omega'_1,\ldots,\omega'_\ell\}\subset\Fq\setminus\Omega\).  A compiled witness is a matrix
\[
W\in\Fq^{n\times s}.
\]
For each row \(i\), the prover interpolates a polynomial \(P_i\in\Fq[X]\) of degree \(<s+\ell\) such that
\(P_i(\omega_j)=W_{i,j}\) for all \(\omega_j\in\Omega\), while the values on \(\Omega'\) are random masks.  The
committed witness is therefore a family of row polynomials \(P=(P_1,\ldots,P_n)\).

SmallWood proves families of constraints over these rows.  Parallel constraints are slotwise identities; aggregated
constraints are domain-summed identities.  In ARC-SPRUCE, the target-shaped commitment equation and the signature
equation are linear coefficient-wise constraints; the inverse-witness equation
\((B_3-\mathbf 1_nx_1)\Had Z=1_n\) is bilinear; and the PRF checks contain the nonlinear Poseidon2-style
S-box constraints.  Range, carry, packing-consistency, and norm gadgets may additionally use aggregated checks.
The whole showing relation is therefore not linear, but its only rational-hash nonlinearity is the bilinear inverse
check.

\subsection{Issuance as a SmallWood statement}
\label{subsec:smallwood-issuance}
\label{subsec:presign-compiled}
The public issuance statement is
\[
  x^{\mathrm{SW}}_{\mathrm{Issue}}=(C_M,A_s,c),
\]
and the hidden witness is
\[
  w^{\mathrm{SW}}_{\mathrm{Issue}}=(M,\vecm,\veck,s,e).
\]
The committed witness rows encode the semantic message \(M\), the message/key split \((\vecm,\veck)\), the
commitment randomness \(s\in\R^{k_s}\), and the error \(e\in\R^{n_c}\).  Semantically, the proof certifies:
\begin{align*}
  c&=C_MM+A_s s+e,\\
  M&=\vecm\Vert\veck,\\
  \veck&\in\mathcal K_{\mathsf{key}},
\end{align*}
together with coefficient bounds on \(M,s,e\) and any issuance policy constraints on \(\vecm\).  The issuance proof
carries no rational-hash inverse constraint and no target computation constraint.  The issuer samples the rational-hash
values only after this proof verifies.

\subsection{Showing as a SmallWood statement}
\label{subsec:smallwood-showing}
\label{subsec:showing-compiled}
The public showing statement is
\[
  x^{\mathrm{SW}}_{\mathrm{Show}}=(\nonce,\prftag,\mA,B,C_M,A_s,\BoundMsg,\BoundSig),
\]
with verifier state handled outside the proof.  The issuance commitment \(c\) is not included in the public showing
statement.  The hidden witness is
\[
  w^{\mathrm{SW}}_{\mathrm{Show}}=(\vecu,M,\vecm,\veck,s,e,\muSig,x_0,x_1,Z).
\]
The committed rows encode the short signature witness \(\vecu\), the semantic message and split \((M,\vecm,
\veck)\), the commitment opening \((s,e)\), the DKLW rational-hash characteristic and randomness
\((\muSig,x_0,x_1)\), the inverse witness \(Z\), and the PRF companion rows.

Semantically, the proof certifies the following coefficient-domain statement:
\begin{align}
  (B_3-\mathbf 1_nx_1)\Had Z&=1_n, \label{eq:sw-show-inv}\\
  \mA\vecu&=B_0+B_1\muSig+B_2x_0+Z+C_MM+A_s s+e, \label{eq:sw-show-sig}\\
  M&=\vecm\Vert\veck, \label{eq:sw-show-encoding}\\
  \prftag&=\PRF(\veck,\nonce), \label{eq:sw-show-prf}
\end{align}
together with key-space membership and the coefficient bounds on all bounded witness components.  Eliminating
\(Z\) yields
\begin{equation}
  (B_3-\mathbf 1_nx_1)\Had
  \bigl(\mA\vecu-B_0-B_1\muSig-B_2x_0-C_MM-A_s s-e\bigr)=1_n,
  \label{eq:sw-show-eliminated}
\end{equation}
which is the coefficient-domain residual used in security arguments.

\subsection{Replay / transform implementation}
\label{subsec:replay-implementation}
Some implementations check the showing relation in a transform or replay basis.  Let
\[
  \mathcal F_{\mathrm{rep}}:\Rq\to\mathcal T
\]
be the public replay map, with block projections \(\mathcal F_{\mathrm{rep},b}\).  The prover then carries replay rows
for the images of
\[
  \mA\vecu,
  \quad
  \muSig,
  \quad
  x_0,
  \quad
  x_1,
  \quad
  C_MM+A_s s+e,
  \quad
  Z,
\]
and exact bridge constraints tie these rows to the coefficient-domain rows.  A replay residual has the form
\begin{equation}
  (B_{3,b}^{\Theta}-\widehat x_{1,b})\Had
  \bigl(\widehat T_b-B_{0,b}^{\Theta}-B_{1,b}^{\Theta}\Had \widehat\mu_{\mathsf{sig},b}
  -B_{2,b}^{\Theta}\Had\widehat x_{0,b}-\widehat c_b\bigr)
  =\mathcal F_{\mathrm{rep},b}(1),
  \label{eq:sw-replay-residual}
\end{equation}
where \(\widehat T_b\) is the replay image of \(\mA\vecu\), \(\widehat c_b\) is the replay image of
\(C_MM+A_s s+e\), and \(B_{i,b}^{\Theta}\) denotes the public replay image of the relevant BB-tran coefficient block.

\begin{assumption}[Exact replay map]\label{ass:replay-transform}
The public replay map \(\mathcal F_{\mathrm{rep}}\) used by the showing implementation is additive, multiplicative on
all products used in the compiled residual, and injective on the coefficient-domain values that appear in the showing
relation.  If a concrete replay representation is not injective, the compiled relation must include explicit linking
constraints sufficient to recover the same coefficient-domain identities.
\end{assumption}

\begin{theorem}[Replay-space certification implies the coefficient-domain committed-message identity]
\label{thm:replay-to-coeff}
Assume Assumption~\ref{ass:replay-transform}.  Assume further that the compiled showing relation enforces the
exact bridges above and the replay residual \eqref{eq:sw-replay-residual}.  Then the extracted coefficient-domain
witness satisfies
\[
  (B_3-\mathbf 1_nx_1)\Had
  \bigl(\mA\vecu-B_0-B_1\muSig-B_2x_0-C_MM-A_s s-e\bigr)=1_n
\]
in \(\Rq\).
\end{theorem}

\begin{proof}
Substitute the exact bridge equalities into the replay residual.  Additivity and multiplicativity of
\(\mathcal F_{\mathrm{rep}}\) transform the replay residual into the replay image of the coefficient-domain residual
\[
  (B_3-\mathbf 1_nx_1)\Had
  \bigl(\mA\vecu-B_0-B_1\muSig-B_2x_0-C_MM-A_s s-e\bigr)-1_n.
\]
Injectivity then forces this coefficient-domain residual itself to vanish, which is exactly the displayed identity.
\end{proof}

\begin{corollary}[Abstract committed-message witness consequence]\label{cor:showing-to-abstract}
Under the assumptions of Theorem~\ref{thm:replay-to-coeff}, the extracted witness additionally satisfies
\[
  (B_3-\mathbf 1_nx_1)\Had Z=1_n,
  \qquad
  \mA\vecu=h_{\muSig,\chiSig}(B)+C_MM+A_s s+e,
  \qquad
  \prftag=\PRF(\veck,\nonce).
\]
\end{corollary}

\begin{proof}
The inverse equation \((B_3-\mathbf 1_nx_1)\Had Z=1_n\) is part of the compiled relation.  By
Theorem~\ref{thm:replay-to-coeff}, the bracketed residual multiplied by \(B_3-\mathbf 1_nx_1\) is \(1_n\).  Combining
this with the inverse equation identifies the bracketed term with \(Z\), hence the displayed signature equation.  The
PRF equation is already part of the showing relation.
\end{proof}

\subsection{PRF companion relation}
\label{subsec:prf-checkpointed}
The public tag is tied to the same hidden key carried in \(M=\vecm\Vert\veck\) through a companion PRF relation.
Selector constraints bind the PRF key lanes to the key rows \(\veck\) and enforce
\(\veck\in\mathcal K_{\mathsf{key}}\); checkpoint constraints certify the nonlinear Poseidon2-style S-box steps;
linear layers and round constants are replayed from public parameters; and the final feed-forward and truncation
constraints bind the public pair \((\nonce,\prftag)\).

A crucial invariant is that the PRF key rows are not separate from the signature/message rows: the compiled showing
statement is a single committed witness family.  If auxiliary alias rows are introduced for engineering convenience,
they must be tied back to the original key rows by explicit copy constraints.  This prevents a prover from using one
hidden key for the committed-message signature relation and another hidden key for the PRF relation.

\subsection{Statement strength}
\label{subsec:statement-strength}
The SmallWood proof hides a family of row polynomials, not merely a tuple of scalars.  Random masks on \(\Omega'\)
statistically hide unopened row values, while the DECS/LVCS/PCS stack reveals only the evaluations needed to check
the PACS constraints.  Consequently, the verifier learns only the public statement: during issuance this is \(c\)
together with public parameters, and during showing it is \((\nonce,\prftag)\) together with the public verification
context.  The verifier does not see \(c\), \(M\), \(s\), \(e\), \(\muSig\), \(x_0\), \(x_1\), \(Z\), or \(\vecu\) in the
showing protocol.

\begin{corollary}[Consequence of an accepting showing proof]\label{cor:showing-proof-consequence}
If the showing NIZK is knowledge sound for the compiled showing relation, then every accepting proof yields a
coefficient-domain witness satisfying \eqref{eq:sw-show-inv}--\eqref{eq:sw-show-prf} and the stated bounds.
\end{corollary}
